\chapter*{\runtitle}

\noindent
The \emph{Community Membership Hiding} (CMH) problem asks: given a graph, a community
detection algorithm, and a target node $u$, find a minimal set of edge rewirings
adjacent to $u$ such that the algorithm no longer assigns it to its original community,
while preserving the overall structural integrity of the graph.
This problem arises in the context of network privacy, where detection algorithms can
inadvertently reveal sensitive affiliations---political or ideological---solely from
a node's structural position.

Recent work formulated CMH as a Markov Decision Process (MDP) solved via Deep
Reinforcement Learning (DRL) using an Actor-Critic (A2C) architecture.
In this thesis we identify a critical architectural limitation of that approach: the
agent evaluates the graph globally without explicitly conditioning on the target node,
thereby learning a generic community-shuffling policy rather than a targeted hiding
strategy.

We propose a novel DRL architecture that explicitly incorporates the target node as an
input feature, employs Graph Attention Networks v2 (GATv2) as the graph encoder to
leverage a more expressive, dynamic attention mechanism, and combines static and
real-time structural signals that track topological changes during rewiring.
We extend the single-agent system to a Dual-Agent Architecture that specializes separate
agents for edge addition and deletion, and introduce a Policy-Guided Subgraphing
heuristic that makes the framework tractable on large graphs.

Extensive experimental evaluation on the standard benchmark datasets demonstrates that
our architecture consistently outperforms both the original DRL agent and strong
gradient-based baselines, achieving high hiding success rates with minimal structural
perturbation.

\bigskip

\noindent\textbf{Keywords:} community detection, community membership hiding, deep
reinforcement learning, graph attention networks, network privacy, GATv2, Actor-Critic.
